\documentclass{letter}

% Top margin, right margin, left margin, bottom margin, footnote skip
\usepackage[tmargin=3.5cm,rmargin=3cm,lmargin=3cm,bmargin=3cm]{geometry} 
\usepackage[utf8]{inputenc}
\usepackage{biblatex}
\addbibresource{./references.bib}
% linktocpage shall be added to snippets.
\usepackage{hyperref,theoremref}

\usepackage{setspace}
\usepackage{enumitem}

\linespread{1.2}

\address{2 Winchester Rd \\ Oxford}
\name{Me}
\signature{My Signature line}

\title{Application to Stipendiary Lectureship in Computer Science}
\author{Yuhao Han} 
\date{\today}

\begin{document}
% \tableofcontents

\begin{letter}{Jesus College \\ University of Oxford}
\opening{Dear Members of the Selection Committee,}
I am writing to apply for Stipendiary Lectureship in Computer Science in Jesus College.
I am currently completing my master's study in Mathematics and Foundation of Computer Science at Oxford and will begin my DPhil study in quantum computing in Michaelmas Term 2026.

This position appears appealing to me, as it would allow me to contribute to the social and intellectual life of the college. 
I am also passionate about teaching, and I am excited to share my knowledge and grow together with my students through small group tutorials settings.

I have a strong educational background that spans computer science, programming, and mathematics.
During my undergraduate and master's studies, I have consistently achieved excellent marks of above 80.
I have research experiences on the Linux kernel, compilers, AI, and algebraic geometry. 
I also have developed substantial software and am familiar with several programming languages including Python, C++, and Rust.
My mathematics training includes linear algebra, homological algebra, category theory, topology, among others.
These experiences qualified me to teach various undergraduate courses in computer science, such as algorithms, imperative and functional programming, and linear algebra, among others. I am also willing to undertake any other duties associated with this role, including providing pastoral care, marking, and helping in admission and college events.

Since my undergraduate studies, I have engaged in various teaching and mentoring activities. 
At the University of Edinburgh, I have volunteered as a tutor to teach first year students linear algebra, and as a mentor to organise activities and provide pastoral care for junior students.
I have worked as a English tutor in China, teaching both adults and high achieving school pupils.
Through these experiences I have gained important perspectives on the academic and psychological needs of students of different backgrounds and confidence levels.
They have also help me develop my communication skills and the ability to explain complex ideals in accessible ways.

Thank you very much for considering my application.
Professor Aleks Kissinger at aleks.kissinger@cs.ox.ac.uk and Professor Sergii Strelchuk at sergii.strelchuk@cs.ox.ac.uk have agreed to be my referees.
I am looking forward to further discuss my application with you in the future.

Yours faithfully,\\
Yuhao Han

\newpage

\textbf{Teaching Statement}

Throughout my teaching career, I have been adhered to the following two teaching principles
\begin{enumerate}
   \item Teaching is more effective with ample examples, practice, and exercises.
   \item Each student has unique talents and perspective, and teaching shall adjust to student's needs
\end{enumerate}
During teaching, I will present students with ample examples and exercise. I would encourage them to complete exercises and guide them to understanding if they are stuck.
I will not simply present to them answers or let them to simple memorise the facts.
I will also actively ask students for their feedbacks and adjust my teaching methods accordingly.

For the subject of computer science, in particular, I believe practice is as important, if not more important, as the theories.
Therefore, I would encourage the students to implement the theories they learned into actual software on their computer.
I will also present to the students practical tools in programming, such as Git, Bash, and Ssh, encourage them to try different ways of programming, such as using command line and Linux, and teach them industrial standards in software engineering.

In recent years, generative AI has posted both opportunities and challenges in education. 
I will follow the University's and the College's statements on AI.
I would encourage students to use AI to aid their study when appropriate and also remind them to prioritise their own skill in coding, writing, and understanding complex concepts.
I would in particular warn them not to rely on AI for studying, as AI' hallucination (making up of fake facts) could be extremely detrimental.

At last, I believe teaching is a mutually beneficial activity, as framed by the Confucius saying \textit{study and teaching complements each other}.
 as it exposes me to diverse perspectives and helps me gain better understanding of the subject.


\textbf{Research Interests}

My interests in mathematics and computer science stem from the following firm belief: the mathematics, while itself beautiful and attractive enough, when applied correctly, can create great positive impacts on our society through computer science.
The prime evidence of this is the internet and artificial intelligence, both  have their foundation in advanced mathematics.
For my current master's dissertation, I am exploring a new tensor contraction algorithm combining graph theory and ZX calculus, which is a categorical representation of linear maps into visually intuitive graphs.
I am actively exploring the mathematics theory of contraction algorithms, as well as implementation of the algorithm using Rust.

My DPhil project will be a theoretical investigation of quantum circuit complexity using geometric methods. It studies how quantum operations can be modelled as paths on spaces of unitary transformations, with the aim of understanding the resources needed to implement these operations using standard
quantum gates. 
The research will analyze reduced and symmetry-constrained models, including few-qubit systems and quotients by permutation symmetries or the Clifford group, in order to make the mathematical problem more tractable. 
The work is purely mathematical and computational, using differential geometry, Lie groups, Clifford algebra, quantum information theory, and
symbolic/numerical computation. 

These research project are directly related to undergraduate teaching, as they involves advance ideas in linear algebra, algorithms, and proof systems. 
I would also introduce my research to students to help them understand the frontier of researches.

\end{letter}
\end{document}
